% Included by appendix-voltage-response.tex.
\subsection{Choosing voltages for specified addition-energy changes}
\label{app:virtual-gates}
\label{sec:operating-points-virtual-gates}

Loading requires moving an addition energy below a reservoir chemical
potential; a compensated exchange pulse instead holds the addition
thresholds fixed while changing tunneling. Because one electrode can
produce both effects, we use the combined response to all varied
voltages to select the desired change. For this local calculation,
the response can be measured without recalculating the full device
potential for every pulse.

The many-electron Hamiltonian identifies what this measured
response represents.
For a small voltage change, write
\[
 \delta H=\sum_j \widehat W_j\,\delta V_j,\qquad
 \widehat W_j=\left.\frac{\partial H}{\partial V_j}\right|_{\mathbf V^0}.
\]
For an isolated, nondegenerate state on a specified charge branch,
the Hellmann--Feynman relation gives
\[
 \frac{\partial(E_N-E_{N-1})}{\partial V_j}
 =\langle\widehat W_j\rangle_N-\langle\widehat W_j\rangle_{N-1}.
\]
This is the addition-energy derivative in
Eq.~\eqref{eq:device-addition-response}, expressed through the
voltage-dependent Hamiltonian.

Let $\mu_i(\mathbf V)$ be the addition energy of dot $i$ for a
specified charge transition. Near reference voltages $\mathbf V^0$,
the response to the $N_g$ varied gate electrodes is
\begin{equation}
 d\mu_i=\sum_{j=1}^{N_g}G_{ij}\,dV_j,\qquad
 G_{ij}=\left.\frac{\partial\mu_i}{\partial V_j}\right|_{\mathbf V^0},
 \qquad d\boldsymbol\mu=G\,d\mathbf V.
 \label{eq:calibration-response-matrix}
\end{equation}
The index $j$ includes plunger and barrier electrodes.
The reservoir contact voltages and magnetic field are held fixed
in the examples below. Together with the chosen charge
configuration, these reference settings define the
\emph{operating point}. Each matrix entry has units of energy
per voltage. For a finite small step, the differential relation
gives $\Delta\boldsymbol\mu\simeq G\Delta\mathbf V$.

Suppose we want to increase $\mu_1-\mu_2$ without changing
$(\mu_1+\mu_2)/2$. The required energy changes have opposite signs,
whereas a pulse on one plunger may shift both energies in the
same direction. We use the measured response matrix to find
voltages whose combined effects give the desired pair of changes.

Consider two dots controlled by two plunger voltages $V_1,V_2$
and one barrier voltage $V_{\mathrm B}$. The full relation is
\begin{equation}
 \begin{pmatrix}d\mu_1\\d\mu_2\end{pmatrix}
 =
 \underbrace{\begin{pmatrix}
 G_{11}&G_{12}&G_{1\mathrm B}\\
 G_{21}&G_{22}&G_{2\mathrm B}
 \end{pmatrix}}_{G}
 \begin{pmatrix}dV_1\\dV_2\\dV_{\mathrm B}\end{pmatrix}.
 \label{eq:calibration-full-gate-response}
\end{equation}

\subsubsection{Obtaining the response matrix}
The control calculation needs all six entries, including those of
the barrier electrode. The charge-transition slopes and
finite-bias measurements provide these entries;
we now assemble them into $G$. On a dot--reservoir transition,
$\mu_i=\mu_{\mathrm R}$. Vary electrode $j$ and adjust the
corresponding plunger $V_i$ to stay on that transition, with the
remaining electrode fixed. The measured slope is
\[
 c_{ij}=\left.\frac{dV_i}{dV_j}\right|_{\mu_i=\mu_{\mathrm R}},
 \qquad j\ne i.
\]
Since $0=G_{ii}dV_i+G_{ij}dV_j$, it gives
$G_{ij}=-G_{ii}c_{ij}$. The four slopes
$c_{12},c_{1\mathrm B},c_{21},c_{2\mathrm B}$ determine the
cross-responses relative to the two plunger responses.
Measure the latter,
$g_1=G_{11}$ and $g_2=G_{22}$, using the finite-bias boundaries
of Section~\ref{sec:measurement-bias-spectroscopy}.
The complete matrix is then
\begin{equation}
 G=\begin{pmatrix}
 g_1&-g_1c_{12}&-g_1c_{1\mathrm B}\\
 -g_2c_{21}&g_2&-g_2c_{2\mathrm B}
 \end{pmatrix}.
 \label{eq:calibration-matrix-from-lines}
\end{equation}
The transition slopes give ratios; the finite-bias measurement
supplies the energy scale. Use nearby smooth segments for the
specified charge sectors, since a change in a neighbouring
occupation can shift a transition discontinuously.

\subsubsection{Solving for the voltage change}
With $G$ known, the requested change $d\boldsymbol\mu^\star$ is the
input and $d\mathbf V$ is the unknown. The control problem is
$G\,d\mathbf V=d\boldsymbol\mu^\star$.
A rank-two $2\times3$ matrix leaves one free voltage direction:
two energy conditions do not uniquely determine three voltages.
To write the solutions explicitly, partition the already measured
matrix as
\[
 G=\begin{pmatrix}G_{\mathrm P}&\mathbf g_{\mathrm B}\end{pmatrix},
 \qquad
 d\mathbf V_{\mathrm P}=\begin{pmatrix}dV_1\\dV_2\end{pmatrix},
\]
where $G_{\mathrm P}$ contains the first two columns and
$\mathbf g_{\mathrm B}$ is the third. If $G_{\mathrm P}$ is
invertible, specifying $dV_{\mathrm B}$ gives
\begin{equation}
 d\mathbf V_{\mathrm P}
 =G_{\mathrm P}^{-1}
       \left(d\boldsymbol\mu^\star-\mathbf g_{\mathrm B}dV_{\mathrm B}\right).
 \label{eq:calibration-virtual-gates}
\end{equation}
Thus keeping the barrier voltage fixed is one possible choice;
the full response matrix also permits solutions that change it.
Coordinated voltage directions chosen to control specified
physical quantities are called \emph{virtual gates}.

\subsubsection{An opposite shift of two addition energies}
For a numerical illustration, take
\[
 G=\begin{pmatrix}-100&-25&-10\\-20&-80&-15\end{pmatrix}
       \frac{\mu\mathrm{eV}}{\mathrm{mV}}.
\]
These are illustrative response values, not experimental data.
Let $\varepsilon_\mu=\mu_1-\mu_2$ and
$\bar\mu=(\mu_1+\mu_2)/2$.
Increasing $\varepsilon_\mu$ by $80\,\mu\mathrm{eV}$ at fixed
$\bar\mu$ requires
$(\Delta\mu_1,\Delta\mu_2)=(40,-40)\,\mu\mathrm{eV}$.
Choosing $\Delta V_{\mathrm B}=0$ gives
\begin{equation}
 \begin{pmatrix}\Delta V_1\\\Delta V_2\\\Delta V_{\mathrm B}\end{pmatrix}
 \simeq
 \begin{pmatrix}-0.56\\0.64\\0\end{pmatrix}\mathrm{mV}.
 \label{eq:calibration-two-gate-example}
\end{equation}
Substitution gives $56-16=40$ and $11.2-51.2=-40$,
in $\mu\mathrm{eV}$, for the two energy changes.
More generally, with $dV_{\mathrm B}=0$,
\begin{equation}
 d\mathbf V_{\mathrm P}
 =G_{\mathrm P}^{-1}
   \begin{pmatrix}1&1/2\\1&-1/2\end{pmatrix}
   \begin{pmatrix}d\bar\mu\\d\varepsilon_\mu\end{pmatrix}.
 \label{eq:calibration-common-differential}
\end{equation}
The difference $\varepsilon_\mu$ concerns the chosen addition
transitions. Its relation to the orbital detuning
$\epsilon_1-\epsilon_2$ also depends on interaction energies.

The calculation has now supplied a voltage vector for the requested
energy shift. Apply that small step, then locate the transitions
again. Their displacements, converted to energy with the measured
responses, test whether the requested energy changes occurred.
If the chosen $G_{\mathrm P}$ is nearly singular, its inverse
requires large compensating voltages. Another independent pair
of columns of the full $G$ may provide a better set of controls.

\subsection{Changing tunneling at fixed addition energies}

An exchange pulse should strengthen the interdot coupling without
unintentionally moving the dot--reservoir loading thresholds.
We therefore seek a voltage direction that keeps both addition
energies fixed but changes $t_c$. To first order, the first
requirement is simply
\begin{equation}
 G\,d\mathbf V=0.
 \label{eq:calibration-fixed-energy-nullspace}
\end{equation}
For the rank-two matrix above, its nullspace is one-dimensional.
If $G_{\mathrm P}$ is invertible, choose a vector in that
nullspace with barrier component equal to one:
\begin{equation}
 \mathbf n=
 \begin{pmatrix}-G_{\mathrm P}^{-1}\mathbf g_{\mathrm B}\\1\end{pmatrix},
 \qquad
 d\mathbf V=\mathbf n\,dv,\qquad G\mathbf n=0.
 \label{eq:calibration-null-direction}
\end{equation}
Here $dv=dV_{\mathrm B}$ parametrizes the coordinated voltage
change. Equivalently, after choosing the barrier step, the
plunger compensation is
\begin{equation}
 d\mathbf V_{\mathrm P}
 =-G_{\mathrm P}^{-1}\mathbf g_{\mathrm B}\,dV_{\mathrm B}.
 \label{eq:calibration-barrier-compensation}
\end{equation}
The barrier pulse and its two compensating plunger pulses are
applied together. Their energy shifts cancel even though the
potential between the dots, and hence the wave-function overlap,
can change.

For the illustrative matrix above,
\[
 \mathbf n=
 \begin{pmatrix}-17/300\\-13/75\\1\end{pmatrix}.
\]
A $0.30\,\mathrm{mV}$ barrier step therefore requires plunger
steps $-0.017\,\mathrm{mV}$ and $-0.052\,\mathrm{mV}$.
Their combined first-order shifts of both $\mu_1$ and $\mu_2$
vanish. Whether this direction changes tunneling is a separate,
measurable property. Define
$w_j=\partial t_c/\partial V_j$ at the operating point.
Along the coordinated direction,
\begin{equation}
 dt_c=\sum_j w_j\,dV_j
     =\left(\sum_j w_j n_j\right)dv.
 \label{eq:calibration-tunneling-direction}
\end{equation}
A nonzero coefficient permits tunneling control at fixed addition
energies. If it vanishes, these three controls cannot achieve
that change to first order at this operating point.

To measure the coefficient, apply several small values of $v$
in $\mathbf V(v)=\mathbf V^0+\mathbf n v$.
At each setting use the charge-distribution measurement of
Section~\ref{sec:measurement-equilibrium-charge} to determine
$t_c(v)$, and recheck the addition thresholds.
The response must be measured along this coordinated path:
changing the plungers also affects tunneling.

The scan gives $t_c(v)$ along a path that keeps the addition
energies fixed. For spin gates, the same voltage path can be
characterized by its exchange coupling $J(v)$ using
Section~\ref{sec:measurement-exchange-oscillations}.

For several couplings, append response rows such as
$d\ln(t_{ij}/t_{\mathrm{ref}})
=\sum_g B_{ij,g}\,dV_g$, where $t_{\mathrm{ref}}>0$ is a fixed
reference energy. The combined matrix can prescribe a change
of $t_{12}$ with $t_{23}$ and the addition energies held fixed,
provided its rank permits those independent constraints.
Such control of energies and tunnel couplings has been
demonstrated using virtual gates~\cite{Hsiao2020}.

\subsubsection{Turning voltage settings into a pulse}
The energy-shifting and compensated barrier directions solve different
problems: the first loads or unloads electrons, whereas the second
changes their coupling without moving the loading thresholds. A change
between two settings means a change of the physical voltage vector,
not an instantaneous change of a model parameter by itself.
For a locally linear path, a pulse can be programmed as
\[
 \mathbf V(t)=\mathbf V^{\mathrm A}
       +f(t)\bigl(\mathbf V^{\mathrm B}-\mathbf V^{\mathrm A}\bigr),
 \qquad 0\le f(t)\le1.
\]
The function $f$ rises from zero to one, holds, and eventually
returns or proceeds to another setting. For an exchange pulse,
the displacement can follow the compensated direction
$\mathbf n$; for loading, it must instead move an addition energy
across a reservoir chemical potential. The endpoint conditions
depend on whether the operation loads an electron, retains a fixed
charge sector, or converts a spin state into charge motion.

The compensation is local. For a larger excursion, measure or
update the response along the path. The static matrix sets the
voltage combinations; the measured time-dependent response sets
the pulse durations and tests the effect of the edges.
